\chapter{USA}

\section{The Constitutional Foundation}

\section{The 10th Amendment --- The Root of Everything}

The US Constitution (ratified 1788) does \textbf{not} mention education anywhere. The 10th Amendment states:

\begin{important}
	``The powers not delegated to the United States by the Constitution, nor prohibited by it to the States, are reserved to the States respectively, or to the people.''
\end{important}

This one sentence is why the USA has \textbf{no national curriculum, no national textbook board, and no national textbook.} Education is constitutionally a \textbf{state power}, not a federal power.

This is the complete opposite of Bangladesh, where NCTB is a \textbf{national, centralized} authority.

\section{There Are 50 Separate Education Systems}

The USA has \textbf{50 states}. Each state legally operates its own education system. Technically, there is no single ``American education system''~--- there are \textbf{50 systems} running in parallel, with some federal overlay (explained in Section 2).

%%─────────────────────────────────────────────────────────────────────────────
\section{The Federal Government's Role}

\section{The US Department of Education (DoED)}

\begin{itemize}
	\item Created by Congress in 1979 via the Department of Education Organization Act.
	\item Does \textbf{NOT} write curriculum.
	\item Does \textbf{NOT} mandate textbooks.
	\item Does \textbf{NOT} set graduation requirements.
	\item Does \textbf{NOT} run any school.
\end{itemize}

\section{What the Federal Government Actually Does}

\begin{itemize}
	\item Distribute federal \textbf{funding} to states (via formulas and grants).
	\item Enforce federal \textbf{civil rights laws} in schools (\eg no discrimination).
	\item Collect education data and statistics nationally.
	\item Run federal student loan programs for higher education.
	\item Set conditions tied to federal money (see Section 3).
\end{itemize}

\section{Current Status of the DoED (July 2026)}

Under the Trump administration (returned January 2025), the DoED has been significantly restructured. An executive order was signed targeting its dismantlement. Many functions were transferred to other departments (\eg student loans to Treasury, special education to Health and Human Services). As of mid-2026, the DoED exists in a drastically reduced form. \textbf{This is an active, ongoing situation --- not yet fully resolved by legislation.}

Key point: Even when the DoED existed at full power, it never had the authority NCTB has. Its dismantlement makes curriculum even more decentralized.

%%─────────────────────────────────────────────────────────────────────────────
\section{State-Level Authority}

\section{States Set the Standards}

Each of the 50 states has a \textbf{State Board of Education} and a \textbf{State Department of Education}. These bodies:
\begin{itemize}
	\item Write or adopt \textbf{academic content standards}.
	\item Set \textbf{graduation requirements}.
	\item Set \textbf{teacher certification requirements}.
	\item Determine \textbf{school year length}.
	\item Handle \textbf{state funding distribution} to local districts.
\end{itemize}

\section{Academic Standards --- What Are They?}

Standards are \textbf{not} textbooks. They are documents that say, for example: ``By the end of 4th grade, students should be able to multiply two-digit numbers.'' Standards tell you \textbf{what} to teach. They do not tell you \textbf{which book} to use to teach it.

\section{The Common Core Experiment (2010--present)}

In 2010, the National Governors Association and the Council of Chief State School Officers (non-government bodies) created the \textbf{Common Core State Standards (CCSS)} for Math and English Language Arts.

\begin{itemize}
	\item These were \textbf{voluntary} --- no law forced states to adopt them.
	\item The federal government incentivized adoption via the \textbf{Race to the Top} grant program.
	\item At peak adoption: \textbf{45 states + DC} adopted Common Core.
	\item As of 2023--2024: Several states \textbf{withdrew} (Indiana 2014, Oklahoma 2014, South Carolina 2015, Texas/Alaska/Virginia never adopted).
	\item So Common Core is \textbf{not a national curriculum} --- it was a voluntary alignment attempt that partially fragmented.
\end{itemize}

\section{State Curriculum Authority}

Research shows: in approximately \textbf{43 states}, curriculum selection authority rests primarily with \textbf{local school districts}, not the state government. This means: even within one state, different school districts can teach using different textbooks and curricula.

%%─────────────────────────────────────────────────────────────────────────────
\section{Local School Districts}

\section{What Is a School District?}

A \textbf{school district} is a local educational agency (LEA) --- a government body that directly runs public schools in a defined geographic area.

The USA has approximately \textbf{13,000 school districts} (US Census Bureau data).

Each district has a locally elected or appointed \textbf{School Board} that decides:
\begin{itemize}
	\item Which textbooks to purchase.
	\item Which curriculum programs to implement.
	\item Local school policies and tax levies.
\end{itemize}

\section{Implication: Textbook Chaos (Or Freedom)}

In districts with local curriculum authority, a 7th-grade math teacher in one Texas district might use a completely different textbook than a teacher in a neighboring district. There is no law against this.

%%─────────────────────────────────────────────────────────────────────────────
\section{The Textbook System in America}

\section{Adoption States vs.\ Open Territory States}

\textbf{Adoption States ($\sim$19--20 states):} The state government maintains an \textbf{official approved list} of textbooks for K-12 subjects. Districts must select from that approved list. Examples: Texas, California, Florida, North Carolina, Virginia.

\textbf{Open Territory States ($\sim$30--31 states):} No state-approved list. Local districts choose any textbook from any publisher.

\section{Why Texas and California Are Disproportionately Powerful}

Texas ($\sim$5.4M K-12 students) and California ($\sim$6M K-12 students) are the two largest state markets. Because they are adoption states with large budgets, textbook publishers \textbf{design national editions to satisfy Texas and California first}. Then sell the same books to the rest of the country.

Texas and California effectively shape what's in American textbooks nationwide, even though they have no legal authority over other states.

\section{The Private Textbook Publishing Industry}

Unlike Bangladesh (where NCTB prints and distributes government textbooks for free for Class 1--10), the USA has a \textbf{private, for-profit textbook publishing industry}. Major publishers: Pearson, McGraw-Hill, Houghton Mifflin Harcourt (HMH), Bedford Freeman \& Worth, Cengage. These are corporations. Schools and districts \textbf{purchase} textbooks with public (tax) funds.

\section{No Free Textbooks in Most of K-12}

In public K-12, students use textbooks owned by the school and return them at year end. Students don't usually buy or keep them --- different from Bangladesh where students own their NCTB books.


%%─────────────────────────────────────────────────────────────────────────────
\section{The K-12 Grade Structure}

\section{Basic Structure}

\begin{center}\small
	\begin{tabularx}{0.9\textwidth}{@{} l l X @{}}
		\toprule
		\textbf{Level}            & \textbf{Grades} & \textbf{Approximate Ages} \\
		\midrule
		Kindergarten              & K               & 5--6                      \\
		Elementary School         & 1--5 or 1--6    & 6--11                     \\
		Middle School/Junior High & 6--8 or 7--8    & 11--14                    \\
		High School               & 9--12           & 14--18                    \\
		\bottomrule
	\end{tabularx}
\end{center}


``K-12'' = Kindergarten through Grade 12. Approximately 13 years of schooling before university.

\section{Variation Exists}

The exact split varies by state and district. Some places run K-8 then 9-12 high school. There is no single mandated structure nationally.

\section{No National School Leaving Exam}

There is \textbf{no equivalent to SSC or HSC} at the national level. Each state may have its own \textbf{state assessment tests} (STAAR in Texas, MCAS in Massachusetts, PARCC in some states). These vary enormously.

Historically, students took the \textbf{SAT} (College Board) or \textbf{ACT} (ACT Inc.) for college admissions --- but these are \textbf{private standardized tests}, not government exams, and many universities have now made them optional.

%%─────────────────────────────────────────────────────────────────────────────
\section{American Higher Education}

\section{Structure of Higher Education}


\begin{center}\small
	\begin{tabularx}{0.97\textwidth}{@{} >{\raggedright\arraybackslash}p{3.5cm} X @{}}
		\toprule
		\textbf{Type}         & \textbf{Description}                                                  \\
		\midrule
		Community College     & 2-year public colleges, offer Associate degrees, low cost             \\
		Liberal Arts College  & 4-year private colleges, small, focus on undergraduate education      \\
		State University      & Public 4-year universities, funded by state government                \\
		Research University   & Public or private, offer Bachelor + Master + PhD, major research      \\
		For-Profit University & Private companies operating universities (controversial, many closed) \\
		\bottomrule
	\end{tabularx}
\end{center}


\section{No National University}

The USA has approximately \textbf{5,000+ degree-granting institutions}. There is \textbf{no single institution} analogous to Bangladesh's National University. Each American university grants its \textbf{own degree} in its own name. No central body grants the actual degree.

\section{No National University Entrance Exam}

There is \textbf{no equivalent to Bangladesh's admission exam system}. Each university runs its own \textbf{individual admissions process} evaluating: high school GPA and transcript, SAT/ACT scores (if required), essays, extracurriculars, and letters of recommendation.

Thousands of universities = thousands of separate applications and processes.

\section{No National Higher Education Curriculum}

Universities choose their own degree programs, textbooks (professors choose individually), and assessment methods. A professor teaching Introductory Statistics at University of Michigan can assign a completely different textbook than a professor at UCLA.

\textbf{University textbooks are expensive private market products.} A single textbook can cost \$100--\$300+ USD. This is a well-documented crisis in American higher education.

%%─────────────────────────────────────────────────────────────────────────────
\section{Accreditation}

\section{What Replaces Government Certification of Universities?}

Since the government doesn't run universities, quality is controlled through \textbf{accreditation} --- a peer-review process conducted by private non-profit bodies.

\section{The Seven Regional Accrediting Bodies}


\begin{center}\small
	\begin{tabularx}{0.85\textwidth}{@{} l X @{}}
		\toprule
		\textbf{Body} & \textbf{Region}               \\
		\midrule
		NECHE         & New England                   \\
		MSCHE         & Middle States                 \\
		HLC           & North Central (large area)    \\
		SACSCOC       & Southern states               \\
		WSCUC         & Western (CA, HI, Pacific)     \\
		NWCCU         & Northwest                     \\
		ACCJC         & California community colleges \\
		\bottomrule
	\end{tabularx}
\end{center}


A university \textbf{must} be regionally accredited for its degrees to be widely recognized and for students to use federal financial aid. Unaccredited institutions' degrees are generally not recognized by employers or other universities.

\section{Accreditation Is Not Government-Run}

These bodies are \textbf{private non-profits}, not government agencies. The government recognizes them but does not operate them. This is a fundamentally American model --- private self-regulation with government recognition.


%%─────────────────────────────────────────────────────────────────────────────
\clearpage
\partbanner{PART TWO: EUROPEAN COUNTRIES}


%%─────────────────────────────────────────────────────────────────────────────

\chapter{EHEA}

\section{What Is It?}

The \textbf{Bologna Process} is an intergovernmental agreement (not an EU initiative) that began with the Bologna Declaration signed in 1999 by 29 European education ministers.

Goal: Make higher education degrees \textbf{comparable and compatible} across Europe, so a degree from Poland is recognizable in Spain, and students can transfer credits across borders.

\section{The European Higher Education Area (EHEA)}

The EHEA is the zone created by the Bologna Process.
\begin{itemize}
	\item \textbf{49 countries} are members (as of 2024): all 27 EU members + non-EU countries like UK, Switzerland, Turkey, Serbia, Ukraine, \etc (Russia suspended 2022).
	\item This is \textbf{larger than the EU}.
\end{itemize}

\section{The Three-Cycle Degree System}


\begin{center}\small
	\begin{tabularx}{0.85\textwidth}{@{} l l X @{}}
		\toprule
		\textbf{Cycle} & \textbf{Degree} & \textbf{Typical Duration}  \\
		\midrule
		First          & Bachelor's      & 3 years (180 ECTS credits) \\
		Second         & Master's        & 2 years (120 ECTS credits) \\
		Third          & Doctorate/PhD   & 3+ years                   \\
		\bottomrule
	\end{tabularx}
\end{center}

\textbf{ECTS} = European Credit Transfer and Accumulation System. 1 ECTS $\approx$ 25--30 hours of student workload.


\section{What Bologna Does NOT Do}

Bologna does \textbf{not} harmonize: course content or curriculum, textbooks, admission standards, grading systems, or tuition fees. Each country still controls its own university content. Bologna only standardizes the \textbf{structure and duration} of degrees, and enables credit transfer.


%%─────────────────────────────────────────────────────────────────────────────

\chapter{Finland}


\section{Basic structure}

Finland's education is controlled at the national level through the Finnish National Agency for Education (EDUFI), which writes the \textbf{National Core Curriculum}. Local municipalities and schools then develop their own curricula based on this framework.

Key facts:
\begin{itemize}
	\item No national standardized tests until the Matriculation Examination (\textit{Ylioppilastutkinto}) at end of upper secondary.
	\item Teachers have high autonomy in choosing teaching materials.
	\item No high-stakes testing during K-9.
	\item Consistently high PISA rankings (OECD).
\end{itemize}


\vspace{1em}
\noindent\textcolor{Rule}{\rule{\linewidth}{0.4pt}}
\vspace{0.5em}





%%─────────────────────────────────────────────────────────────────────────────

\chapter{Norway}


\section{Basic structure}

Finland's education is controlled at the national level through the Finnish National Agency for Education (EDUFI), which writes the \textbf{National Core Curriculum}. Local municipalities and schools then develop their own curricula based on this framework.

Key facts:
\begin{itemize}
	\item No national standardized tests until the Matriculation Examination (\textit{Ylioppilastutkinto}) at end of upper secondary.
	\item Teachers have high autonomy in choosing teaching materials.
	\item No high-stakes testing during K-9.
	\item Consistently high PISA rankings (OECD).
\end{itemize}



%%─────────────────────────────────────────────────────────────────────────────



\chapter{South Korea}


\section{Basic structure}

Finland's education is controlled at the national level through the Finnish National Agency for Education (EDUFI), which writes the \textbf{National Core Curriculum}. Local municipalities and schools then develop their own curricula based on this framework.

Key facts:
\begin{itemize}
	\item No national standardized tests until the Matriculation Examination (\textit{Ylioppilastutkinto}) at end of upper secondary.
	\item Teachers have high autonomy in choosing teaching materials.
	\item No high-stakes testing during K-9.
	\item Consistently high PISA rankings (OECD).
\end{itemize}



%%─────────────────────────────────────────────────────────────────────────────



\chapter{Japan}


\section{Basic structure}

Finland's education is controlled at the national level through the Finnish National Agency for Education (EDUFI), which writes the \textbf{National Core Curriculum}. Local municipalities and schools then develop their own curricula based on this framework.

Key facts:
\begin{itemize}
	\item No national standardized tests until the Matriculation Examination (\textit{Ylioppilastutkinto}) at end of upper secondary.
	\item Teachers have high autonomy in choosing teaching materials.
	\item No high-stakes testing during K-9.
	\item Consistently high PISA rankings (OECD).
\end{itemize}



%%─────────────────────────────────────────────────────────────────────────────


\chapter{Singapore}


\section{Basic structure}

Finland's education is controlled at the national level through the Finnish National Agency for Education (EDUFI), which writes the \textbf{National Core Curriculum}. Local municipalities and schools then develop their own curricula based on this framework.

Key facts:
\begin{itemize}
	\item No national standardized tests until the Matriculation Examination (\textit{Ylioppilastutkinto}) at end of upper secondary.
	\item Teachers have high autonomy in choosing teaching materials.
	\item No high-stakes testing during K-9.
	\item Consistently high PISA rankings (OECD).
\end{itemize}



%%─────────────────────────────────────────────────────────────────────────────


\chapter{Netherlands}


\section{Basic structure}

Finland's education is controlled at the national level through the Finnish National Agency for Education (EDUFI), which writes the \textbf{National Core Curriculum}. Local municipalities and schools then develop their own curricula based on this framework.

Key facts:
\begin{itemize}
	\item No national standardized tests until the Matriculation Examination (\textit{Ylioppilastutkinto}) at end of upper secondary.
	\item Teachers have high autonomy in choosing teaching materials.
	\item No high-stakes testing during K-9.
	\item Consistently high PISA rankings (OECD).
\end{itemize}



%%─────────────────────────────────────────────────────────────────────────────

\chapter{Switzerland}


\section{Basic structure}

Finland's education is controlled at the national level through the Finnish National Agency for Education (EDUFI), which writes the \textbf{National Core Curriculum}. Local municipalities and schools then develop their own curricula based on this framework.

Key facts:
\begin{itemize}
	\item No national standardized tests until the Matriculation Examination (\textit{Ylioppilastutkinto}) at end of upper secondary.
	\item Teachers have high autonomy in choosing teaching materials.
	\item No high-stakes testing during K-9.
	\item Consistently high PISA rankings (OECD).
\end{itemize}



%%─────────────────────────────────────────────────────────────────────────────


\chapter{Key Takeaways}


\textbf{1. The USA is the most decentralized major education system in the developed world.} There is no NCTB equivalent. No national textbook. No national exit exam. 50 states, $\sim$13,000 districts, all making independent decisions. The federal government's role is funding + civil rights enforcement, not curriculum.

\textbf{2. Europe is not one system.} The EU legally cannot touch education curriculum (Article 165 TFEU). Each of the 27 EU member states runs its own system. The Bologna Process aligns only the \textit{structure} of university degrees (3+2+3 years), not content.

\textbf{3. Germany = decentralized within Europe} (16 states, each sovereign over education --- similar logic to USA states).

\textbf{4. France = centralized within Europe} (national curriculum, national exam, national civil servant teachers --- closest to Bangladesh's model among major Western countries).

\textbf{5. Bangladesh's NCTB model for K-10 is actually relatively rare} among the world's major economies. The centralized national textbook + free distribution model is more common in developing countries. Developed countries tend toward either centralized curriculum with market textbooks (France) or decentralized everything (USA, Germany).


\source{US Constitution 10th Amendment; Department of Education Organization Act 1979; TFEU Article 165; Bologna Declaration 1999; KMK official documentation; French Ministry of National Education; NCTB Bangladesh official mandate. Data verified through 2025--2026.}
ENDOFFILE
